\chapter{Software Architecture}

\section{Technology Stack}

The web application is built on a deliberately lean stack, chosen to cover all functional
requirements without unnecessary complexity.

\begin{description}[leftmargin=2cm, style=nextline, itemsep=4pt]
  \item[Python / Flask]
    Flask is a lightweight Python micro-framework providing both a REST API and, via the
    \code{flask-socketio} extension, real-time WebSocket communication. The availability
    of a first-class Python binding for Snap7 (\code{python-snap7}) was a decisive factor
    in choosing Python.

  \item[Socket.IO]
    Rather than the browser polling an endpoint every few seconds, Socket.IO maintains a
    persistent WebSocket connection. The server \textit{pushes} new data to all connected
    clients as soon as the poll loop fires. This eliminates redundant requests and allows
    multiple browser windows to display the same live data simultaneously.

  \item[Docker]
    Containerisation provides a reproducible execution environment. The same Docker image
    runs on the development PC (Windows, Docker Desktop), on a Raspberry Pi (Linux ARM),
    and on the SIMATIC Unified panel (Linux x86\_64 via Industrial Edge). Dependencies are
    locked in \code{requirements.txt} and baked into the image at build time.

  \item[Chart.js]
    A JavaScript charting library that renders responsive time-series charts directly in
    the browser without any server-side processing. The identification page uses a dual-Y-axis
    configuration with \code{animation: false} for smooth streaming updates.
\end{description}

Table~\ref{tab:techstack} summarises the technology choices.

\begin{table}[H]
  \centering
  \caption{Technology choices and justifications.}
  \label{tab:techstack}
  \begin{tabular}{@{}L{3.2cm}L{3.5cm}L{6.3cm}@{}}
    \toprule
    \textbf{Layer} & \textbf{Technology} & \textbf{Justification} \\
    \midrule
    Web server       & Flask (Python 3.11) & Lightweight, REST + SocketIO, Snap7 binding \\
    Real-time push   & Socket.IO 5         & WebSocket with HTTP long-poll fallback, multi-client \\
    PLC comm.\ (direct) & Snap7 (TCP 102) & Open-source S7 library, absolute DB access \\
    PLC comm.\ (tags)   & Open Pipe       & Native WinCC Unified tag interface \\
    Frontend charts  & Chart.js 4          & Client-side rendering, no server overhead \\
    Frontend layout  & Bootstrap 5         & Responsive grid, touch-friendly components \\
    Containerisation & Docker / Compose    & Cross-platform (PC $\rightarrow$ Edge panel) \\
    Data persistence & CSV (append-only)   & Simple, portable, no database dependency \\
    \bottomrule
  \end{tabular}
\end{table}

\section{Application Structure}

\subsection{Module Overview}

The application is structured as a set of cooperating Python modules, each responsible
for a single concern:

\begin{table}[H]
  \centering
  \caption{Python module responsibilities.}
  \label{tab:modules}
  \begin{tabular}{@{}L{4.5cm}L{8.5cm}@{}}
    \toprule
    \textbf{Module} & \textbf{Responsibility} \\
    \midrule
    \code{main.py}                  & Flask routes, Socket.IO events, background thread management \\
    \code{plc\_connector.py}        & Snap7 TCP connection to DB1 \\
    \code{openpipe\_connector.py}   & Open Pipe Unix socket, 21 WinCC tags \\
    \code{cascade\_controller.py}   & PI+P inner loop, P outer loop, mode state machine \\
    \code{identification.py}        & Step test orchestration, 63.2\% FOPTD identification, IMC gains \\
    \code{thermal\_model.py}        & Discrete-time 1st-order simulation fallback \\
    \code{logger.py}                & Append-only CSV writer \\
    \bottomrule
  \end{tabular}
\end{table}

Each class is instantiated once at application startup (\textit{singleton pattern}),
ensuring that all HTTP route handlers and background threads share the same state objects.

\subsection{Threading Model}

The application runs four concurrent threads (Table~\ref{tab:threads}).

\begin{table}[H]
  \centering
  \caption{Background thread responsibilities.}
  \label{tab:threads}
  \begin{tabular}{@{}L{3.0cm}C{1.5cm}L{8.5cm}@{}}
    \toprule
    \textbf{Thread name} & \textbf{Period} & \textbf{Role} \\
    \midrule
    Flask / SocketIO (main) & event    & HTTP request handling, WebSocket dispatch \\
    \code{sim-loop}         & 3\,s     & Advance thermal simulation model (fallback only) \\
    \code{plc-poll}         & 3\,s     & Read DB1 via Snap7, emit \code{process\_data} \\
    \code{openpipe-poll}    & 3\,s     & Read 21 tags, run cascade/identification, emit events \\
    \bottomrule
  \end{tabular}
\end{table}

\textbf{Thread safety.}
Shared control objects (cascade controller, step identifier) are written only from the
\code{openpipe-poll} thread. HTTP route handlers only set parameters and flags that are
read at the next poll cycle. This single-writer design avoids race conditions without
requiring explicit locks on the control objects.

\section{PLC Communication Layer}

\subsection{Snap7 --- Direct DB Access}

\code{plc\_connector.py} wraps the \code{python-snap7} library. The S7-1500 connection
uses rack~0 and slot~1 (mandatory for all S7-1500 CPUs). Reading DB1 uses absolute byte
offsets:

\begin{lstlisting}[language=Python, caption={Snap7 DB1 read in \texttt{plc\_connector.py}.}]
raw = client.db_read(db_number=1, start=0, size=20)
temperature = snap7.util.get_real(raw,  0)   # DBD0
flow_rate   = snap7.util.get_real(raw,  4)   # DBD4
setpoint_T  = snap7.util.get_real(raw,  8)   # DBD8
setpoint_F  = snap7.util.get_real(raw, 12)   # DBD12
valve_state = snap7.util.get_int( raw, 16)   # DBW16
\end{lstlisting}

Connection failures are handled gracefully: if Snap7 cannot connect (PLC off, wrong IP),
the connector sets \code{\_connected = False} and the poll loop uses the simulation
fallback. There is no application crash.

\subsection{Open Pipe --- WinCC Tag Access}

\code{openpipe\_connector.py} uses Python's standard \code{socket} library to communicate
with the Unix domain socket. The key method \code{read\_all\_tags()} sends a single JSON
request listing all 21 tag names and receives a single JSON response with all values.
This batched read avoids the latency of 21 individual round-trips.

\begin{lstlisting}[language=Python,
  caption={Open Pipe batch read in \texttt{openpipe\_connector.py}.}]
def _read_via_socket(self, tag_names: list) -> dict:
    request = json.dumps({"action": "read", "tags": tag_names}) + "\n"
    with socket.socket(socket.AF_UNIX, socket.SOCK_STREAM) as s:
        s.settimeout(2.0)
        s.connect(self._sock_path)
        s.sendall(request.encode())
        response = b""
        while True:
            chunk = s.recv(4096)
            if not chunk or b"\n" in response:
                break
            response += chunk
    return json.loads(response.split(b"\n")[0].decode())
\end{lstlisting}

Writing is done tag by tag via \code{write\_tag(name, value)}, since control outputs
($F_{1,SP}$) are written one at a time.

\section{Real-Time Data Flow}

\subsection{Poll Loop Logic}

The \code{openpipe-poll} thread executes the following sequence every 3\,s:

\begin{enumerate}[leftmargin=*, itemsep=3pt]
  \item Read all 21 tags via \code{openpipe.read\_all\_tags()}
  \item If identification is running: call \code{ident.feed()} $\rightarrow$ may return
    a new $F_{1,SP}$ to write (identification has \textbf{priority} over cascade)
  \item Otherwise, if cascade mode $\neq$ MANUAL: call \code{cascade.update()} $\rightarrow$
    returns new $F_{1,SP}$; write to PLC via \code{openpipe.write\_tag('F1\_SP', sp)}
  \item Emit \code{process\_tags} Socket.IO event to all connected clients
  \item Emit \code{cascade\_data} event with controller state
  \item If identification active: emit \code{ident\_update} with progress and chart data
\end{enumerate}

\subsection{Socket.IO Events}

\begin{table}[H]
  \centering
  \caption{Socket.IO events emitted by the server.}
  \label{tab:socketio}
  \begin{tabular}{@{}L{3.0cm}L{3.0cm}L{7.0cm}@{}}
    \toprule
    \textbf{Event} & \textbf{Direction} & \textbf{Data payload} \\
    \midrule
    \code{process\_data}  & server $\to$ client & 5 DB1 variables + source indicator \\
    \code{process\_tags}  & server $\to$ client & All 21 WinCC tags + source indicator \\
    \code{cascade\_data}  & server $\to$ client & Mode, setpoints, errors, $F_{1,SP}$ output \\
    \code{ident\_update}  & server $\to$ client & Phase, progress \%, chart data point, results \\
    \bottomrule
  \end{tabular}
\end{table}

Each page subscribes only to the events it needs: the Synoptique listens to
\code{process\_tags}; the Cascade page listens to all four; the Monitor page listens to
\code{process\_data}.

\section{REST API}

Table~\ref{tab:api} summarises the HTTP endpoints.

\begin{longtable}{@{}C{1.5cm}L{5.5cm}L{6.0cm}@{}}
  \caption{REST API endpoint summary.}
  \label{tab:api}\\
  \toprule
  \textbf{Method} & \textbf{Route} & \textbf{Description} \\
  \midrule
  \endfirsthead
  \multicolumn{3}{l}{\small\itshape (continued)}\\
  \toprule
  \textbf{Method} & \textbf{Route} & \textbf{Description} \\
  \midrule
  \endhead
  \bottomrule
  \endfoot
  GET  & \code{/}                        & Monitor page (HTML) \\
  GET  & \code{/schema}                  & Synoptique P\&ID page (HTML) \\
  GET  & \code{/cascade}                 & Cascade identification / control page (HTML) \\
  GET  & \code{/process}                 & Process tag monitor (HTML) \\
  GET  & \code{/api/data}                & Current 5-variable Snap7 snapshot (JSON) \\
  POST & \code{/api/setpoint}            & Update temperature or flow setpoint (DB1) \\
  GET  & \code{/api/process/data}        & All 21 WinCC tags via Open Pipe (JSON) \\
  POST & \code{/api/process/write}       & Write \code{F1\_SP} or \code{F2\_SP} \\
  GET  & \code{/api/cascade/status}      & Cascade controller state (JSON) \\
  POST & \code{/api/cascade/mode}        & Switch MANUAL / AUTO-INNER / AUTO-FULL \\
  POST & \code{/api/cascade/setpoint}    & Set $T_{out2,SP}$ or $T_{in1,SP}$ \\
  POST & \code{/api/cascade/params}      & Set inner $K_p$, $T_i$ and outer $K_{p,ext}$ \\
  POST & \code{/api/identification/start}  & Launch step test \\
  POST & \code{/api/identification/cancel} & Abort; restore $F_{1,SP}$ to baseline \\
  GET  & \code{/api/identification/status} & State, progress, results (JSON) \\
  GET  & \code{/api/export/csv}          & Download full process log (CSV) \\
  GET  & \code{/health}                  & Docker healthcheck (\texttt{\{"status":"ok"\}}) \\
\end{longtable}

\section{Data Persistence}

\code{logger.py} implements an append-only CSV writer. Each poll cycle a new row is
written to \code{/app/logs/process\_data.csv} containing timestamp, all key process
variables, and the \code{source} field.

The log directory is bound to a Docker volume (\code{./logs:/app/logs}), ensuring
survival across container restarts and image rebuilds. The endpoint
\code{GET /api/export/csv} streams the file with a timestamped filename
(\code{process\_data\_YYYYMMDD\_HHMMSS.csv}).

\section{Web Interface Pages}

\subsection{Synoptique (\texttt{/schema}) --- Panel Display}

The Synoptique is a full-screen SVG P\&ID diagram optimised for the
SIMATIC Unified 7\textquotedblright{} touchscreen (800~$\times$~480 pixels). CSS-animated
dashed strokes create visual flow indicators (orange for the hot circuit, blue for the
cold circuit). Ten live values are displayed as overlay cards updated every 3\,s via
\code{process\_tags} Socket.IO events. Touch-sensitive increment/decrement buttons
adjust $F_{1,SP}$ and $F_{2,SP}$ in 0.1\,\si{\metre\cubed\per\hour} steps.

\subsection{Cascade (\texttt{/cascade}) --- Engineering Tool}

Designed for PC use during commissioning. Two tabs:

\begin{itemize}[leftmargin=*, itemsep=3pt]
  \item \textbf{Identification:} step test configuration, live Chart.js dual-axis chart
    ($T_{in1}$ and $T_{out2}$ on left axis; $F_{1,SP}$ step on right axis), progress bar,
    and IMC results table with ``Apply \& Switch'' buttons.
  \item \textbf{Cascade Control:} visual block diagram, parameter forms for inner
    ($K_p$, $T_i$) and outer ($K_{p,ext}$) gains, mode buttons (MANUAL / AUTO-INNER /
    AUTO-FULL), and live status grid showing setpoints, measurements, and errors.
\end{itemize}

\subsection{Process (\texttt{/process}) --- Debug Monitor}

Displays all 21 WinCC tags in three sections (Furnace, HE1 Hot Side, HE1 Cold Side).
Valve tags show colour-coded badges: green~=~OPEN, yellow~=~PARTIAL, red~=~CLOSED.

\subsection{Monitor (\texttt{/}) --- Historical Chart}

Scrolling Chart.js time-series of temperature and flow rate from the Snap7 poll (DB1
values). Includes basic alarm logic, setpoint controls, and a CSV export button.
